\begin{statement}{B.35}{Proved}{Exact first-stage signed normalizers and the surviving midpoint character}
Use the actual sources of B.31, put $\Delta=C/F'$, and let
$A=\widehat{\mathbb Z}^{\times}$.  Write
\[
 \eta_8(\mu)=\frac{\mu_2^2-1}{8}\pmod2,
 \qquad
 \mathcal E_2=\{(t,\mu)\in\widehat{\mathbb Z}\rtimes A:
                         t\bmod2=\eta_8(\mu)\},
\]
where $(t,\mu)(v,\nu)=(t+\mu v,\mu\nu)$.
Then the images of the two actual signed normalizers in
$\operatorname{Aut}(\Delta)$ are
\[
\begin{aligned}
 \operatorname{im}N_\rho^*
   &=\{\operatorname{Int}(\rho^j)h_\mu:
                       j\in\mathbb Z/3,\ \mu\in A\}
     \simeq C_3\rtimes_{\epsilon_3}A,\\
 \operatorname{im}N_\sigma^*
   &=\{\operatorname{Int}((xy)^t\sigma^e)s_\mu:
              e\in\mathbb Z/2,\ (t,\mu)\in\mathcal E_2\}
     \simeq C_2\times\mathcal E_2.
                                                        \tag{B.35.1}
\end{aligned}
\]
Here the normalizer action is the specified strict action of B.31,
all displayed coordinates are unique, and
\[
 \begin{aligned}
 v_3(\mu)&=(1-\epsilon_3(\mu))/2,&
 a_\mu&=(\epsilon_3(\mu)\mu-1)/6,\\
 h_\mu&=\operatorname{Int}((xy^{-1})^{a_\mu}
                         \sigma^{v_3(\mu)})s_\mu.
 \end{aligned}
                                                        \tag{B.35.2}
\]
The quotient $\operatorname{im}N_\sigma^*\to A$ does not split.
In particular the retained $\sigma$ conjugator carries a genuine
additional coordinate: its multiplier-one translation image is
$2\widehat{\mathbb Z}(x+y)$, although the full quotient centralizer
contains $\widehat{\mathbb Z}(x+y)$.

The images of $G_D$, $\mathcal M_D$, and $S_0\mathcal M_D$,
using their projections in B.31, are all
\[
                  \operatorname{Inn}(\Delta)\rtimes s_A,
       \qquad Z(\Delta)=\langle\omega\rangle.             \tag{B.35.3}
\]
If $u|_\Delta=\operatorname{Int}(e_0^{r(u)})s_{\mu(u)}$,
the action kernel of $G_D$ is the concrete set
\[
 \{[c,u]:\mu(u)=1,\quad
                \overline c\,e_0^{r(u)}\in\langle\omega\rangle\}.
                                                        \tag{B.35.4}
\]
The two larger fibre-product kernels are its inverse images.
For the normalizers their kernels are the inverse images of
$(j,\mu)=(0,1)$ and $(e,t,\mu)=(0,0,1)$ in (B.35.1).
These conclusions concern the first modular stage $C/F'$;
they do not collapse the higher Fermat actions to this stage.
\end{statement}
\begin{proof}
We first identify the centralizers which are needed both for
normalization and for the geometric comparison.  Theorems 1 and 2
of \citet[printed p.~31]{B35HerfortRibes1984} state that torsion
elements of a free profinite product are conjugate into a factor,
and that the centralizer of a nontrivial factor element lies in
that factor.  The class of all finite groups satisfies their
closure hypotheses.  The opened primary text is their four-page
announcement with proof outline; we use its stated theorems,
not a discrete reduced-word argument applied to profinite words.
Applying the centralizer theorem to
$C/\langle\omega\rangle=C_3\widehat *C_2$ and lifting gives
\[
 C_C(\rho)=\langle\rho\rangle,
 \qquad C_C(\sigma)=\langle\sigma\rangle.                 \tag{B.35.5}
\]
The finite quotient $H_1=\operatorname{Dic}_3$ has corresponding
centralizers $\langle\overline\rho\rangle$ and
$\langle\overline\sigma\rangle$.  Strip these finite components
in $\Delta$ and use the matrices $R,S$ of B.33.  Since
$M^R=0$ and $M^S=\widehat{\mathbb Z}(x+y)$, this proves
\[
 C_\Delta(\rho)=\langle\rho\rangle,
 \qquad
 C_\Delta(\sigma)=\langle\sigma\rangle
                         \times\widehat{\mathbb Z}(x+y).
                                                        \tag{B.35.6}
\]
The kernel of $H_1\to\operatorname{Aut}(M)$ is its central
$\langle\omega\rangle$, and $M^{H_1}=0$.  Thus a central
element of $\Delta$ is exactly $1$ or $\omega$.

By B.33 the image of $G_D$ is generated by
$\operatorname{Inn}(\Delta)$ and $s_A$.  These intersect trivially:
on $M$, an inner automorphism acts by an element of the displayed
$S_3$ representation, while $s_\mu$ acts by the scalar $\mu$;
the only scalar in this integral $S_3$ representation is the identity.
Thus $\mu=1$ and $s_\mu=1$.  This proves the semidirect product
in (B.35.3), and directly gives (B.35.4).  Both normalizers map
onto $\gamma$: for a representative $u\in D$ of a class in
$\gamma$, a conjugator for $u(\rho)$, respectively $u(\sigma)$,
corrects $u$ to an element of the relevant normalizer.
Consequently the projections $\mathcal M_D\to G_D$ and
$S_0\mathcal M_D\to G_D$ are onto.  This proves their image
and kernel assertions without replacing the conjugator source
by a group of unmarked automorphisms.

\emph{The $\rho$ image.}
Put $n=(\mu-1)/2$.  Direct calculation in $\Delta$ gives
\[
 s_\mu(\rho)=y^n\rho,\qquad
 s_\mu(\sigma)=\sigma^{\epsilon_2(\mu)},\qquad
 \sigma\rho\sigma^{-1}=x\rho^{-1}.                       \tag{B.35.7}
\]
For example use $\sigma=e_0e_1e_0$ and
$\rho=e_0^{-1}e_1^{-1}$, together with
$e_0^2=\omega x$, $e_1^2=\omega y$; the central sign in
$s_\mu(\sigma)$ is fixed by its abelianization modulo four.
If $\epsilon_3(\mu)=1$, conjugation by
$(xy^{-1})^{n/3}$ removes the translation in $s_\mu(\rho)$.
If $\epsilon_3(\mu)=-1$, first conjugate by $\sigma$ and then
by $(xy^{-1})^{-(n+1)/3}$.  The divisions are integral because
$\mu$ is odd and has its specified residue modulo three.
This proves $h_\mu(\rho)=\rho^{\epsilon_3(\mu)}$.
Moreover
\[
 a_\mu+\epsilon_3(\mu)\mu a_\nu=a_{\mu\nu}.
\]
The action of $\sigma$ on $x-y$ is $-1$, and its square is
central.  This identity therefore proves
$h_\mu h_\nu=h_{\mu\nu}$ as automorphisms of $\Delta$.
Conjugation by $h_\mu$ acts on $\operatorname{Int}(\rho)$
through $\epsilon_3(\mu)$.

Every multiplier occurs in $N_\rho^*$, by its surjection to
$\gamma$ and B.33.  Any two target normalizers with the same
multiplier differ by $\operatorname{Int}(C_\Delta(\rho))$,
which is exactly $C_3$ by (B.35.6).  This entire finite kernel
already comes from $\langle\rho\rangle\subset C$.
Thus every $h_\mu$ and exactly the stated three elements above
it occur.  This proves the first line of (B.35.1).

\emph{An integral restriction on the $\sigma$ coordinate.}
Equation (B.35.7) and (B.35.6) show that every ambient target
normalizer has a unique expression
$\operatorname{Int}((xy)^t\sigma^e)s_\mu$, with
$t\in\widehat{\mathbb Z}$ and $e\in\mathbb Z/2$.
We determine which are supplied by the actual source.
First suppose its multiplier is one.  Write the source element
as $[c,u]$ and absorb the inner loop factor of B.33 into $c$.
We may therefore assume $u\in D_0$, so that $u$ is the identity
on $\Delta$.  There are unique normalized conjugators
$\alpha,\beta\in F$ satisfying
\[
          u(\rho)=\alpha\rho\alpha^{-1},\qquad
          u(\sigma)=\beta\sigma\beta^{-1}.               \tag{B.35.8}
\]
Existence follows by reducing a conjugator in $H_1$ and
multiplying it by a power of the corresponding finite-order
generator; uniqueness follows from (B.35.5) and
$F\cap\langle\rho\rangle=F\cap\langle\sigma\rangle=1$.
The equations modulo $F'$ give
\[
             \alpha^{\rm ab}=0,\qquad
             \beta^{\rm ab}=k(x+y).
\]
Here $1-R$ is injective over $\widehat{\mathbb Z}$, even
though its determinant three is not a unit at three.

The action of $R$ on $\gamma_2F/\gamma_3F$ is its determinant,
namely one.  Since $\alpha\in F'$, (B.35.8) implies
$u(\rho)\rho^{-1}\in\gamma_3F$.  The fixed loop $\sigma\rho$
then gives $u(\sigma)\sigma^{-1}\in\gamma_3F$ as well.
In the integral class-two quotient write uniquely
\[
             \beta=x^k y^k c^q,\qquad c=[x,y].
\]
The interchange sends this element to
$x^k y^k c^{-k^2-q}$.  Thus
\[
                        2q+k^2=0.
                                                        \tag{B.35.9}
\]
Reduction modulo two forces $k\in2\widehat{\mathbb Z}$.
No division by two in a Malcev algebra has been used.
Finally the normalizer equation says $c\beta\in C_C(\sigma)$,
so its target translation is $-k$.  The actual multiplier-one
image is therefore contained in
$C_2\times2\widehat{\mathbb Z}(x+y)$.

\emph{Midpoint realization in the actual arithmetic subgroup.}
Here we justify the geometric comparison and its basepoint
ambiguity explicitly.  Let $U=\mathbb P^1-\{0,1,\infty\}$,
with tangent $0+1$.  The quotient stack $[U/S_3]$ has orbifold
fundamental group $C_3\widehat *C_2=C/\langle\omega\rangle$:
the invariant map
\[
               j(z)=\frac{256(1-z+z^2)^3}{z^2(1-z)^2}
\]
has stabilizers of orders three and two over $0$ and $1728$,
and its cusp is removed.  Its index-six kernel is $\pi_1(U)$.
In the braid presentation, the two standard meridians are
$a^2,b^2$, whose images are $x^{-1},y^{-1}$.
The standard arithmetic GT action of E.6 restricts to the usual
based Galois action on these meridians.  It is consequently the
arithmetic action on this quotient stack.  Indeed an automorphism
of $C/\langle\omega\rangle$ which fixes $F$ pointwise is the
identity: its discrepancy on any element centralizes $F$.
That centralizer is trivial.  To see this, a centralizing element
has a power in $F$, hence a power in $Z(F)=1$; the torsion and
centralizer theorems quoted above then exclude a nontrivial
finite-order element centralizing the nonabelian open subgroup.
Thus the extension of the specified $F$ action is unique.

The relevant involution of $U$ is $h(z)=1-z$, with rational
fixed point $P=1/2$.  The subgroup above $\langle h\rangle$ is
\[
 F\rtimes\langle\overline\sigma\rangle
     =\langle x,\overline\sigma\rangle
     \simeq\widehat{\mathbb Z}\widehat *C_2:
\]
eliminate $y=\overline\sigma x\overline\sigma^{-1}$ in its
profinite presentation.  The stabilizer at $P$ is an order-two
subgroup of this group.  Herfort--Ribes Theorem 1 conjugates
it to the $C_2$ factor, since $\widehat{\mathbb Z}$ is
torsion-free.  The conjugator can be taken in $F$: multiplying
by its $C_2$ component does not change the conjugated factor.
Choose an etale path from the tangent to $P$ realizing that
conjugacy.  If $g\in G_\mathbb Q$, its path cocycle is a
conjugator $\beta_g$ for the involution, since the stabilizer
at the rational fixed point is constant downstairs.
Upstairs in $C$ the possible extra $\omega$ factor is excluded
by the forced image of $\sigma$ in $C^{\rm ab}$; the exponent
is precisely $\epsilon_2(\chi(g))$.  Hence
$[\beta_g^{-1},g]\in N_\sigma^*$ for the actual arithmetic
source, with the signs of B.31 retained.

The covers obtained by adjoining roots of $z$ and $1-z$
compute the two abelian coordinates of this path cocycle.
Their values at $1/2$ are both the Kummer class of $1/2$;
this is also \citet[Lemma 3.16]{FA_Wickelgren2012}.
The same direct cover calculation is valid at every prime,
including two.  Reversing the path or the two meridian
orientations changes the common sign.  Changing the comparison
path adds $(\chi(g)-1)v$ to its abelian cocycle for some
$v\in\widehat{\mathbb Z}^2$.  Consequently, on
$G_{\mathbb Q(\mu_\infty)}$ its diagonal coordinate is
exactly $\pm\kappa_2(g)$, and for arbitrary $g$ its parity
is that of $\kappa_2(g)$, independently of either convention.
These are the only two assertions needed for the exact image.
The diagonal equality itself also follows from
$g|_\Delta=s_{\chi(g)}$ and (B.35.7).

For completeness the image of this Kummer character on the
cyclotomic kernel is exactly
\[
          \kappa_2(G_{\mathbb Q(\mu_\infty)})
                              =2\widehat{\mathbb Z}.      \tag{B.35.10}
\]
The image is a closed subgroup of $\widehat{\mathbb Z}$.
Its divisibility by $p^r$ is equivalent to
$2^{1/p^r}\in\mathbb Q(\mu_\infty)$.
The square root belongs to $\mathbb Q(\zeta_8)$.
For every integer $m>2$, $X^m-2$ is Eisenstein, and
$\mathbb Q(2^{1/m})$ is real but has nonreal conjugates;
it is therefore not normal over $\mathbb Q$ and cannot
be a finite subextension of the abelian extension
$\mathbb Q(\mu_\infty)/\mathbb Q$.  Thus the two-primary
image has valuation exactly one and every other primary
image has valuation zero, proving (B.35.10).
Finally
$\sqrt2=\zeta_8+\zeta_8^{-1}$ gives
\[
        \kappa_2(g)\pmod2=\eta_8(\chi(g)).                \tag{B.35.11}
\]
Every multiplier occurs arithmetically.  Equations
(B.35.10)--(B.35.11) supply every translation in the claimed
parity coset; (B.35.9), applied to the quotient of any two
actual normalizers having the same multiplier, excludes all
others.  The finite $C_2$ factor already comes from $\sigma$.
This proves the second line of (B.35.1).  Its law is the
stated affine law, because $s_\mu$ scales $x+y$ by $\mu$;
$\operatorname{Int}(\sigma)$ is central in this normalizer
image, since inversion has no effect on its order-two class.

\emph{Nonsplitting and finite certificate.}
A section over $A$ would give a one-cocycle $t:A\to\widehat{\mathbb Z}$
with $t(\mu)\bmod2=\eta_8(\mu)$.  Since $A$ is abelian,
commuting $\mu$ with $-1$ gives
\[
                  2t(\mu)=(1-\mu)t(-1).
\]
But $\eta_8(-1)=0$, so $t(-1)=2b$ and
$t(\mu)=(1-\mu)b$ is always even.  This contradicts any
$\mu$ with $\mu_2\equiv3\pmod8$.  Projection to
$\mathcal E_2$ shows that the extra central $C_2$ cannot
remove the obstruction.

The program \texttt{g\_signed\_normalizers.py} reconstructs
$\Delta_n=C/(F'F^n)$ from the exact integral cocycle of B.24a,
for $1\le n\le8$. It checks (B.35.7), the section law in
(B.35.2), and the reductions of both normalizer images on
all group elements, with all generator-edge homomorphism
checks. It also checks (B.35.9) modulo $2,4,\ldots,64$.
Nonsplitting is already visible with translation modulo four
and multiplier modulo eight: none of the four pairs of lifts
of $3,-1\in(\mathbb Z/8)^\times$ commute.
These are finite checks of the proved source images;
arithmetic realization is supplied by the preceding path
and Kummer argument, not inferred from finite shadows.
\end{proof}
